(1) 通常理想弦的单位长度横向回复力为 \(Ty_{xx}\)；若材料中存在与局部应变率成正比的 Kelvin–Voigt 内耗，则再增加 \(\eta\partial_t y_{xx}\)。方程仍是线性偏微分方程，但存在耗散而不能用保守系统的完备实频率直接代替。质量线密度 \(\mu\)、张力 \(T\) 分别控制惯性与弹性，\(\tau=\eta/T\) 具有时间量纲。
(2) 取位移 \(y=y_0e^{i(kx-\omega t)}\)，有 \(y_{tt}=-\omega^2y\)、\(y_{xx}=-k^2y\)、\(y_{xxt}=i\omega k^2y\)。代入可得 \(\mu\omega^2=k^2(T-i\eta\omega)\)，因此 \(k(\omega)=\omega\sqrt{\mu/(T-i\eta\omega)}\)。选取实部为正、虚部为正的平方根分支，则 \(e^{ikx}\) 在 \(x\) 正方向衰减，符合被动介质传播条件。
(3) 定义 \(c_0=\sqrt{T/\mu}\)，展开 \((1-i\omega\tau)^{-1/2}=1+i\omega\tau/2+O((\omega\tau)^2)\)，得到 \(k\simeq(\omega/c_0)[1+i\omega\tau/2]\)。故振幅衰减长度的倒数 \(\alpha=\operatorname{Im}k\simeq\omega^2\tau/(2c_0)\)，而 \(\operatorname{Re}k=\omega/c_0+O(\tau^2)\)，即相速度没有一阶修正，能量衰减率约为振幅的两倍。
(4) 将原方程乘 \(y_t\) 并在 \([0,L]\) 上积分，两次分部积分且固定端 \(y_t(0,t)=y_t(L,t)=0\)，得到 \(dE/dt=-\eta\int_0^L y_{xt}^2\,dx\le0\)，其中 \(E=\frac12\int_0^L[\mu y_t^2+Ty_x^2]\,dx\)。耗散项是速度梯度平方，严格非负；这在任意波形而不仅在单色波近似下成立。
